%% defense_mechanism.tex -- submission version

\section{Content-Level Defense: \trustorigin{} Tagging}
\label{sec:defense}

\trustorigin{} is a content-level defense for the recovery channel. It
preserves source information at the failure boundary, carries that information
through the recovery pipeline, and selects an untrusted framing for content
whose origin is external or unknown. The mechanism is complementary to
authorization checks at an action sink; it does not make the LLM
instruction-following decision deterministic.

\subsection{Provenance as Data}
\label{sec:trustorigin}

The root failure mode is treating error traces as opaque strings. Once a
string has lost its origin, a recovery strategy cannot distinguish a
framework-generated exception from a web page, file, fixture, or tool return
that happened to appear in the failure. The defense therefore carries the
origin with the value:

\begin{lstlisting}[language=Java,caption={Illustrative provenance-carrying recovery value.},label={lst:tagged-content}]
record TaggedContent(
    String text,
    TrustOrigin origin,
    String sourceRef) {}

enum TrustOrigin {
    SYSTEM_GENERATED,
    TOOL_OUTPUT_INTERNAL,
    TOOL_OUTPUT_EXTERNAL,
    USER_INPUT,
    UNKNOWN
}
\end{lstlisting}

The label is assigned at capture time, not inferred from the text after the
fact. Unknown tools and undeclared sources are treated as external by default.
This fail-closed rule avoids silently granting an internal trust class to a
new tool or plugin.

\subsection{Recovery-Frame Selection}
\label{sec:trustorigin-frame}

The recovery policy uses the origin label to select one of two templates:

\begin{itemize}[leftmargin=*]
    \item \textbf{Verified internal diagnostic:} the model may use the
    diagnostic as framework state, subject to ordinary task and tool policy.
    \item \textbf{External or unknown content:} the content is explicitly
    framed as data to inspect. It is not treated as an instruction, and any
    requested side effect must be independently authorized.
\end{itemize}

The defense is intentionally content-level. A sufficiently instruction-following
model may still obey an injected directive inside an untrusted frame; the
empirical question is therefore the observed reduction under the specified
payloads, not a deterministic security guarantee.

\subsection{Threat Coverage and Scope}
\label{sec:trustorigin-scope}

\trustorigin{} addresses externally originated failure content re-injected
under a recovery frame. It does not authorize role replacement, tool
permission changes, topology mutation, or other privileged side effects. Those
actions require a separate structural check at the orchestrator sink, as
discussed in \S\ref{sec:design-patterns}. The completed defense-aware screen
and its limitations are reported in \S\ref{sec:defense-eval}.
